\documentclass[10pt, letterpaper]{article}

% Packages:
\usepackage[
    ignoreheadfoot, % set margins without considering header and footer
    top=2 cm, % seperation between body and page edge from the top
    bottom=2 cm, % seperation between body and page edge from the bottom
    left=2 cm, % seperation between body and page edge from the left
    right=2 cm, % seperation between body and page edge from the right
    footskip=1.0 cm, % seperation between body and footer
    % showframe % for debugging 
]{geometry} % for adjusting page geometry
\usepackage[explicit]{titlesec} % for customizing section titles
\usepackage{tabularx} % for making tables with fixed width columns
\usepackage{array} % tabularx requires this
\usepackage[dvipsnames]{xcolor} % for coloring text
\definecolor{primaryColor}{RGB}{0, 79, 144} % define primary color
\usepackage{enumitem} % for customizing lists
\usepackage{fontawesome5} % for using icons
\usepackage{amsmath} % for math
\usepackage[
    pdftitle={Alexander Reinthal's CV},
    pdfauthor={Alexander Reinthal},
    pdfcreator={LaTeX with RenderCV},
    colorlinks=true,
    urlcolor=primaryColor
]{hyperref} % for links, metadata and bookmarks
\usepackage[pscoord]{eso-pic} % for floating text on the page
\usepackage{calc} % for calculating lengths
\usepackage{bookmark} % for bookmarks
\usepackage{lastpage} % for getting the total number of pages
\usepackage{changepage} % for one column entries (adjustwidth environment)
\usepackage{paracol} % for two and three column entries
\usepackage{ifthen} % for conditional statements
\usepackage{needspace} % for avoiding page brake right after the section title
\usepackage{iftex} % check if engine is pdflatex, xetex or luatex

% Ensure that generate pdf is machine readable/ATS parsable:
\ifPDFTeX
    \input{glyphtounicode}
    \pdfgentounicode=1
    \usepackage[T1]{fontenc}
    \usepackage[utf8]{inputenc}
    \usepackage{lmodern}
\fi

\usepackage[default, type1]{sourcesanspro} 

% Some settings:
\AtBeginEnvironment{adjustwidth}{\partopsep0pt} % remove space before adjustwidth environment
\pagestyle{empty} % no header or footer
\setcounter{secnumdepth}{0} % no section numbering
\setlength{\parindent}{0pt} % no indentation
\setlength{\topskip}{0pt} % no top skip
\setlength{\columnsep}{0.15cm} % set column seperation
\makeatletter
\let\ps@customFooterStyle\ps@plain % Copy the plain style to customFooterStyle
\patchcmd{\ps@customFooterStyle}{\thepage}{
    \color{gray}\textit{\small Alexander Reinthal - Page \thepage{} of \pageref*{LastPage}}
}{}{} % replace number by desired string
\makeatother
\pagestyle{customFooterStyle}

\titleformat{\section}{
    % avoid page braking right after the section title
    \needspace{4\baselineskip}
    % make the font size of the section title large and color it with the primary color
    \Large\color{primaryColor}
}{
}{
}{
    % print bold title, give 0.15 cm space and draw a line of 0.8 pt thickness
    % from the end of the title to the end of the body
    \textbf{#1}\hspace{0.15cm}\titlerule[0.8pt]\hspace{-0.1cm}
}[] % section title formatting

\titlespacing{\section}{
    % left space:
    -1pt
}{
    % top space:
    0.3 cm
}{
    % bottom space:
    0.2 cm
} % section title spacing

% \renewcommand\labelitemi{$\vcenter{\hbox{\small$\bullet$}}$} % custom bullet points
\newenvironment{highlights}{
    \begin{itemize}[
        topsep=0.10 cm,
        parsep=0.10 cm,
        partopsep=0pt,
        itemsep=0pt,
        leftmargin=0.4 cm + 10pt
    ]
}{
    \end{itemize}
} % new environment for highlights

\newenvironment{highlightsforbulletentries}{
    \begin{itemize}[
        topsep=0.10 cm,
        parsep=0.10 cm,
        partopsep=0pt,
        itemsep=0pt,
        leftmargin=10pt
    ]
}{
    \end{itemize}
} % new environment for highlights for bullet entries


\newenvironment{onecolentry}{
    \begin{adjustwidth}{
        0.2 cm + 0.00001 cm
    }{
        0.2 cm + 0.00001 cm
    }
}{
    \end{adjustwidth}
} % new environment for one column entries

\newenvironment{twocolentry}[2][]{
    \onecolentry
    \def\secondColumn{#2}
    \setcolumnwidth{\fill, 4.5 cm}
    \begin{paracol}{2}
}{
    \switchcolumn \raggedleft \secondColumn
    \end{paracol}
    \endonecolentry
} % new environment for two column entries

\newenvironment{threecolentry}[3][]{
    \onecolentry
    \def\thirdColumn{#3}
    \setcolumnwidth{1 cm, \fill, 4.5 cm}
    \begin{paracol}{3}
    {\raggedright #2} \switchcolumn
}{
    \switchcolumn \raggedleft \thirdColumn
    \end{paracol}
    \endonecolentry
} % new environment for three column entries

\newenvironment{header}{
    \setlength{\topsep}{0pt}\par\kern\topsep\centering\color{primaryColor}\linespread{1.5}
}{
    \par\kern\topsep
} % new environment for the header

\newcommand{\placelastupdatedtext}{% \placetextbox{<horizontal pos>}{<vertical pos>}{<stuff>}
  \AddToShipoutPictureFG*{% Add <stuff> to current page foreground
    \put(
        \LenToUnit{\paperwidth-2 cm-0.2 cm+0.05cm},
        \LenToUnit{\paperheight-1.0 cm}
    ){\vtop{{\null}\makebox[0pt][c]{
        \small\color{gray}\textit{Last updated in March 2025}\hspace{\widthof{Last updated in March 2025}}
    }}}%
  }%
}%

% save the original href command in a new command:
\let\hrefWithoutArrow\href

% new command for external links:
\renewcommand{\href}[2]{\hrefWithoutArrow{#1}{\ifthenelse{\equal{#2}{}}{ }{#2 }\raisebox{.15ex}{\footnotesize \faExternalLink*}}}


\begin{document}
    \newcommand{\AND}{\unskip
        \cleaders\copy\ANDbox\hskip\wd\ANDbox
        \ignorespaces
    }
    \newsavebox\ANDbox
    \sbox\ANDbox{}

    \placelastupdatedtext
    \begin{header}
        \fontsize{30 pt}{30 pt}
        \textbf{Alexander Reinthal}

        \vspace{0.3 cm}

        \normalsize
        \mbox{{\footnotesize\faMapMarker*}\hspace*{0.13cm}Gothenburg, Sweden}%
        \kern 0.25 cm%
        \AND%
        \kern 0.25 cm%
        \mbox{\hrefWithoutArrow{mailto:email@reinthal.me}{{\footnotesize\faEnvelope[regular]}\hspace*{0.13cm}email@reinthal.me}}%
        \kern 0.25 cm%
        \AND%
        \kern 0.25 cm%
        \mbox{\hrefWithoutArrow{tel:+46-73-036-28-78}{{\footnotesize\faPhone*}\hspace*{0.13cm}073-036 28 78}}%
        \kern 0.25 cm%
        \AND%
        \kern 0.25 cm%
        \mbox{\hrefWithoutArrow{https://reinthal.me/}{{\footnotesize\faLink}\hspace*{0.13cm}reinthal.me}}%
        \kern 0.25 cm%
        \AND%
        \kern 0.25 cm%
        \mbox{\hrefWithoutArrow{https://linkedin.com/in/alexander-reinthal}{{\footnotesize\faLinkedinIn}\hspace*{0.13cm}alexander-reinthal}}%
        \kern 0.25 cm%
        \AND%
        \kern 0.25 cm%
        \mbox{\hrefWithoutArrow{https://github.com/reinthal}{{\footnotesize\faGithub}\hspace*{0.13cm}reinthal}}%
    \end{header}

    \vspace{0.3 cm - 0.3 cm}


    \section{Alexander Reinthal}



        
        \begin{onecolentry}
            \textbf{\textit{Senior Data Engineer, Independent Contributor \& Open-Source Advocate}} with expertise in real-time data ingestion, data modeling, and orchestrating cloud-based data workloads. Strong background in \textbf{\textit{Apache Flink, Databricks \& Unity Catalog, and Snowflake ,similar to AWS Glue, S3, and EMR-equivalent architectures}}. Proven open-source contributions to Dagster \& DLT Hub. Passionate about developer experience, mentoring engineers in distributed systems, and ML-driven analytics.
        \end{onecolentry}


    
    \section{Experience}



        
        \begin{twocolentry}{
            Gothenburg, Sweden

        Feb 2022 – present
        }
            \textbf{Knowit Solutions Cocreate}, Data Engineer / Tech Lead / Delivery Manager
            \begin{highlights}
                \item Led an agile R\&D team to build real-time data ingestion pipelines handling massive volumes of CI-pipeline logs using \textbf{\textit{Apache Flink (AWS Kinesis/Data Streams equivalent) \& Spark Structured Streaming}}.
                \item Designed scalable data pipelines on \textbf{\textit{Databricks \& Snowflake, equivalent to AWS Glue \& Athena}}, enabling advanced analytics workloads.
                \item Served as the technical lead, engaging with a client to design and implement a modern data platform leveraging Snowflake, Dagster, and dbt.
                \item Owned the development of an asset-centric data modeling approach for Databricks \& Delta Live Tables, enhancing streaming data reliability and developer experience.
                \item Mentored junior engineers in Python, SQL, DAG-based orchestration, and modern data engineering best practices. Led workshops on aligning technical solutions with business objectives.
                \item Collaborated with stakeholders to refine data governance, mediate upstream data contracts, and improve data visibility using effective metadata management.
            \end{highlights}
        \end{twocolentry}


        \vspace{0.2 cm}

        \begin{twocolentry}{
            Gothenburg, Sweden

        Apr 2019 – Jan 2022
        }
            \textbf{NTT Security}, Security Analyst
            \begin{highlights}
                \item Developed predictive analytics tooling for detecting malicious network traffic, leveraging Python, machine learning, and distributed processing (precursor to modern real-time data pipelines).
                \item Automated security alert processing, optimizing real-time event analysis and reducing response times.
                \item Designed Python-based data ingestion and processing workflows, akin to real-time ETL pipelines in modern data engineering.
                \item Provided mentorship and security training, educating analysts on threat intelligence, anomaly detection, and alert triage automation.
            \end{highlights}
        \end{twocolentry}


        \vspace{0.2 cm}

        \begin{twocolentry}{
            Gothenburg, Sweden

        June 2017 – Oct 2018
        }
            \textbf{Ericsson}, Python Software Engineer
            \begin{highlights}
                \item Led a team to develop a log parsing and analytics tool that scaled into a dedicated engineering team at Ericsson.
                \item Automated data workflow optimizations in Python, improving efficiency in log-based machine-generated data.
            \end{highlights}
        \end{twocolentry}



    
    \section{Projects}



        
        \begin{twocolentry}{
            \href{https://github.com/reinthal/howigotpwned}{GitHub}
        }
            \textbf{HowIgotPwned}
            \begin{highlights}
                \item Developed a modern data lakehouse solution leveraging Apache Iceberg, Flink, Spark, Nessie (similar to AWS Glue), MinIO (AWS S3 equivalent), and Dagster to analyze high-volume data.
                \item Prototyped real-time ingestion pipelines using Apache Flink and Iceberg tables stored on AWS S3, demonstrating cloud-scale flexibility.
            \end{highlights}
        \end{twocolentry}


        \vspace{0.2 cm}

        \begin{twocolentry}{
            \href{https://github.com/reinthal/datalake-stack}{GitHub}
        }
            \textbf{Datalake Stack (On-Prem Equivalent of AWS Data Architecture)}
            \begin{highlights}
                \item Designed and implemented an on-premise data lake solution for Kubernetes using Nessie (similar to AWS Glue), Apache Iceberg, MinIO (AWS S3 equivalent), Apache Flink and Spark.
                \item Integrated modern orchestration frameworks (FluxCD and Traefik) to manage multi-tenant data workloads efficiently.
            \end{highlights}
        \end{twocolentry}


        \vspace{0.2 cm}

        \begin{twocolentry}{
            \href{https://github.com/reinthal}{GitHub}
        }
            \textbf{Open Source Contributions}
            \begin{highlights}
                \item Dagster Contribution: \href{https://github.com/dagster-io/dagster/pull/24188}{PR \#24188} – Improved workflow orchestration, scheduling, and monitoring of data pipelines. 
                \item DLT Hub Contribution: \href{https://github.com/dlt-hub/verified-sources/pull/594}{PR \#594} – Mentored junior engineers, contributing Python data ingestion enhancements.  
            \end{highlights}
        \end{twocolentry}



    
    \section{Technologies}



        
        \begin{onecolentry}
            \textbf{Data Architectures:} Data Modeling, Data Warehousing, Data Lakehouse Solutions
        \end{onecolentry}

        \vspace{0.2 cm}

        \begin{onecolentry}
            \textbf{Semantic Layers:} dbt Cloud, Cube.dev
        \end{onecolentry}

        \vspace{0.2 cm}

        \begin{onecolentry}
            \textbf{Data Catalogs:} \textbf{\textit{Polaris, Snowflake, Unity \& Nessie (Similar to Amazon Glue)}}
        \end{onecolentry}

        \vspace{0.2 cm}

        \begin{onecolentry}
            \textbf{Big Data \& Streaming:} \textbf{\textit{Apache Flink, Apache Spark (AWS EMR equivalent)}}, RabbitMQ
        \end{onecolentry}

        \vspace{0.2 cm}

        \begin{onecolentry}
            \textbf{Data Engineering Tools \& Platforms:} \textbf{\textit{Databricks \& Delta Live Tables, Snowflake (similar to Amazon Athena)}}, dbt, Dagster
        \end{onecolentry}

        \vspace{0.2 cm}

        \begin{onecolentry}
            \textbf{Languages:} Python, SQL, Java, Terraform, Nix, Rust, C++, Assembly, R
        \end{onecolentry}

        \vspace{0.2 cm}

        \begin{onecolentry}
            \textbf{Cloud \& Platforms:} MinIO \& Azure Blob Storage (AWS S3 equivalent), Azure, Kubernetes, Proxmox
        \end{onecolentry}

        \vspace{0.2 cm}

        \begin{onecolentry}
            \textbf{Infrastructure \& DevOps:} Docker, Podman, FluxCD, Git, CI/CD pipelines, Terraform, Helm
        \end{onecolentry}

        \vspace{0.2 cm}

        \begin{onecolentry}
            \textbf{Business Intelligence Tools:} PowerBI, Apache Superset, Lightdash
        \end{onecolentry}


    
    \section{Education}



        
        \begin{threecolentry}{\textbf{MS}}{
            Sept 2016 – June 2018
        }
            \textbf{Chalmers Technical University}, Engineering Physics
            \begin{highlights}
                \item Coursework: Statistical Physics, Neural Networks and Machine Learning
            \end{highlights}
        \end{threecolentry}

        \vspace{0.2 cm}

        \begin{threecolentry}{\textbf{BS}}{
            Sept 2013 – June 2016
        }
            \textbf{Gothenburg University}, Computer Science
            \begin{highlights}
                \item Coursework: Algorithms, Testing, Debugging and Verification, Theoretical Computer Science, Operating Systems, Cryptography, Cyber Security
            \end{highlights}
        \end{threecolentry}


    
    \section{Publications}



        
        \begin{samepage}
            \begin{twocolentry}{
                Nov 2018
            }
                \textbf{Data Modelling for Predicting Exploits}

                \vspace{0.10 cm}

                \mbox{Reinthal, A}, \mbox{Filippakis, E.}, \mbox{Almgren, M.}
                \vspace{0.10 cm}

        \href{https://doi.org/10.1007/978-3-030-03638-6_21}{$10.1007/978-3-030-03638-6_21$}
            \end{twocolentry}
        \end{samepage}


    

\end{document}